% Generated draft: requires mathematical and source review.
\documentclass[11pt]{article}
\usepackage[margin=25mm]{geometry}
\usepackage{fontspec}
\setmainfont{Latin Modern Roman}
\usepackage{amsmath,amssymb,mathtools,unicode-math}
\setmathfont{Latin Modern Math}
\usepackage{enumitem,needspace,xurl,hyperref}
\setlength{\parindent}{0pt}
\setlength{\parskip}{5pt}
\setlength{\emergencystretch}{3em}
\newcommand{\N}{\mathbb{N}}
\newcommand{\Z}{\mathbb{Z}}
\newcommand{\Q}{\mathbb{Q}}
\newcommand{\R}{\mathbb{R}}
\newcommand{\C}{\mathbb{C}}
\newcommand{\F}{\mathbb{F}}
\newcommand{\A}{\mathbb{A}}
\newcommand{\PP}{\mathbb{P}}
\newcommand{\E}{\mathbb{E}}
\newcommand{\eps}{\varepsilon}
\newcommand{\dd}{\,\mathrm{d}}
\begin{document}
\section*{UnsolvedMath ID: 20001898}
\textbf{Source problem code:} catalogue entry
\subsection{Definitions and mathematical statement}
The extracted Question \(6.2.1\), attributed to Burns, asks for a \(C^\infty\), or at least \(C^k\) with \(k\geq3\), metric admitting a nonclosed geodesic along which Gaussian curvature vanishes everywhere. The aid reconstructs the intended setting as a closed surface \(S\) of genus at least two with \(K_g\leq0\). Those hypotheses are missing from the truncated sentence and are essential: without restrictions, flat examples do not represent the intended problem.

With the reconstructed setting explicit, the question is whether there exist such \((S,g)\) and a complete unit-speed geodesic \(\gamma:\R\to S\) for which
\[
 K_g(\gamma(t))=0\quad\text{for every }t\in\R,
\]
but the geodesic-flow orbit is not periodic. Periodic means that for some \(T>0\), both position and velocity repeat after \(T\). The source ends mid-sentence after saying that a negative answer would give a positive answer to another question; the aid identifies that cross-reference as Question \(6.1.1\), whose content is not supplied here.

The aid's local-to-global partial result is as follows. Put \(Z=\{x\in S:K_g(x)=0\}\) and \(C=\overline{\gamma(\R)}\). If there is a neighborhood \(U\) of \(C\) for which \(Z\cap U\) is an embedded one-dimensional submanifold, then every such all-zero-curvature geodesic \(\gamma\) is periodic. This is a condition near the entire base-image closure, not merely a regularity condition along a short arc.

Its second-jet consequence requires \(g\in C^4\), so that \(K_g\in C^2\). On a closed \(C^4\) nonpositively curved surface, every nonperiodic all-zero-curvature geodesic must have
\[
 p\in\overline{\gamma(\R)}\quad\text{such that}\quad
 \operatorname{Hess}_p K_g=0
\]
as a bilinear form. In particular, if
\[
 K_g(p)=0\quad\Longrightarrow\quad
 \operatorname{Hess}_pK_g\ne0
 \qquad\text{for every }p\in S,
\]
then every flat geodesic is closed. The aid includes as a sufficient example a codimension-one Morse--Bott zero-curvature set with negative normal Hessian: its Hessian is nonzero, although it has the tangent direction in its kernel. The condition is that the whole Hessian be nonzero, not that it be nondegenerate.

These restrictions do not resolve the general smooth problem, and the Hessian statement cannot be imposed on an arbitrary \(C^3\) metric whose curvature need only be \(C^1\). They identify necessary degeneracy in any remaining counterexample, not a construction or a nonexistence theorem without the additional hypotheses.
\subsection{Short English statement}
In the reconstructed closed, nonpositively curved genus-at-least-two setting, Burns asks for a nonperiodic geodesic entirely in the zero-curvature locus. The aid proves periodicity when that locus is locally a curve near the orbit closure. For \(C^4\) metrics, any counterexample must accumulate where the entire curvature Hessian vanishes; the general smooth and \(C^3\) questions remain.
\subsection{Sources}
\begin{itemize}
\item Local UnsolvedMath registry, permanent ID \(20001898\), source code \texttt{catalogue entry}; canonical collection \texttt{aim-geometry-notes.json}, notes[235]. Source-item field \(6.2\); actual visible question \(6.2.1\), attributed to Burns.
\item AIM workshop ``Geodesics,'' Burns--Matveev report ``Open problems and questions about geodesics'': \url{https://aimath.org/pastworkshops/geodesicsproblems.pdf}. The canonical sentence is visibly truncated after ``A negative answer to this question would give a positive answer to.'' The aid reconstructs the nonpositive-curvature, closed genus-at-least-two setting and cross-reference to Question \(6.1.1\). No PDF was opened in this session.
\item Supplied research attempt \(1\), \texttt{partial\_result}, problem status at run ``open,'' rigor label ``proved,'' novelty confidence ``low.'' Registry status \texttt{partially\_solved}, classification \texttt{PARTIAL-PROGRESS}, updated 14 August 2026. Original remarks and literature are empty; no full external reference for the obstruction theorem, proposer-field value, or proposed year is supplied.
\end{itemize}
\end{document}
